\specialchapter{记号索引}

SpA(x) （谱）2 R(x,λ) （预解式）3 Ā （单位元的添入）4 Sp′A(x) （谱）4 X(A) （特征标）6 X(h) （特征标函子）6 G\_A(x) （Gelfand 变换）7 X′(A) （特征标）9 X′(h) （特征标函子）9 G′A(x) （Gelfand 变换）10 J(A) （本原理想）12 V(M) （J(A) 的闭集）12 Υ(T) （与 T 相伴的理想）13 cl(π) （π 的类）13 Ā （不可约表示的集合）14 γ （正则表示）16 δ （正则表示）16 B(X;K) （有界函数）17 C\_b(X;K) （有界连续函数）17 C\_0(X;K) （在无穷远处为零的连续函数）17 C(X;K) （连续函数）17 K(X;K) （具有紧支集的连续函数）17 M¹(G) （有界测度）19 ρ(x) （谱半径）21 ev\_x （在 x 处的取值）32 X′ （Alexandroff 紧化）33 Sp\^{}Λ\_A(x) （同时谱）42 O(U;F) （全纯函数）49 O(K;F) （芽）50 O(U;A) （全纯函数代数）50



\(\mathcal{O}(\mathrm{K};\mathrm{A})\) （芽代数）50 \(\Theta_{a}\) （全纯函数演算）51 \(f(a)\) （全纯函数演算的函数记号）72 \(j_{\mathrm{H}}\) （与 H 相伴的幂等元）82 \(\mathrm{R_H}(x,\lambda)\) （关于 \(\mathrm{A_H} = j_{\mathrm{H}}\mathrm{A}j_{\mathrm{H}}\) 的预解式）83 \(\mathcal{O}_{\mathbf{R}}(\mathrm{S})\) （芽的 \(\mathbf{R}\) -代数（满足 \(f = f^{*}\)））86 \(\mathrm{Sp}_{\mathrm{A(C)}}(x)\) （复谱）86 \(x^{*}\) （伴随）96 \(\mathrm{A}_h\) （Hermite 元）96 \(\mathcal{C}'(\mathrm{X})\) （在一点处为零的函数）103 \(f(x)\) （连续函数演算）110 \(\mathrm{A}_{+}\) （正元）115 \(x \geqslant y\) （\(\mathrm{A}_h\) 上的序关系）117 \(x^{+}\) （正部）117 \(x^{-}\) （负部）117 \(|x|\) （绝对值）117 \(x^{\alpha}\) （正元的幂）118 \(\sqrt{x}\) （正元的平方根）118\\
Stell(A) （包络星代数）124\\
Stell(G) （群的星代数）125 \(e_{\mathrm{H}}(u)\) （谱投影）129 \(\mathrm{E_H}(u)\) （谱子空间）129 \(\widetilde{\mathrm{E}}_{\mathrm{H}}(u)\) （谱子空间）129 \(\iota(u)\) （数值值域）135 \(m(u)\) （谱的下确界）138 \(\mathrm{M}(u)\) （谱的上确界）138 \(|u|\) （Hilbert 空间之间态射的绝对值）139 \((j,|u|)\) （极分解）140\\
P(X) （多项式的极限）146 \(\operatorname{lsc}(g)\) （换位子的长度）155 \(\mathbf{C}[\mathrm{G}]\) （群代数）183 \(\operatorname{Stell}_r(\mathrm{G})\) （约化星代数）184 \(\operatorname{Lip}_b(\mathrm{X})\) （有界 Lipschitz 函数）196 \(\widetilde{f}(x)\) （\(\mathrm{L}^1 (\mathrm{G})\) 上的对合）199 \(\widehat{\mathbf{G}}\) （对偶群）201 \(\chi (\mu)\) （\(\mathcal{M}^1 (\mathrm{G})\) 的特征标）202 \(\mathrm{A}^{\perp}\) （正交）205 \(\widehat{\varphi}\) （对偶态射）205 \(\mathcal{F}_{\mathrm{G}}(\mu)\) （Fourier 变换）207 \(\overline{\mathcal{F}}_{\mathrm{G}}(\mu)\) （逆 Fourier 变换）207 \(\mathcal{F}_{\mathrm{G}}(f)\) （Fourier 变换）208 \(\overline{\mathcal{F}}_{\mathrm{G}}(f)\) （逆 Fourier 变换）208\\
A(G) （由 \(f*g\) 生成的子空间）210 \(\widehat{\mathrm{A}}(\mathrm{G})\) （A(G) 在 \(\mathcal{F}\) 下的象）211

B(G) （由 \(f \in L^{1}(G)\)（满足 \(\mathcal{F}(f) \in L^{1}(\widehat{G})\)）组成的子代数）222 T （环面）235 \(\mathcal{S}(\mathbf{R}^{n})\) （Schwartz 函数）238 V(Λ) （余体积）239 \(\mathcal{S}'(\mathbf{R}^{n})\) （缓增广义函数）239 A(W) （W 中的特征标）258 Y(F) （生成的 L∞(G) 的子空间）258 s(t) （符号函数）261 \(\mathcal{P}(X)\) （概率测度）261 ζ(s) （zeta 函数）267 L(s, χ) （Dirichlet L-函数）268 L(G) （G 在 R 中的表示）307 Λ(n) （von Mangoldt 函数）312

